\documentclass[9pt,t]{beamer}
% =====================================================================
% finance-beamer.tex — shared Beamer style for the LLM-in-Finance course
% =====================================================================
% Reusable slide style. Each deck \input's this immediately after
% \documentclass{beamer}. It provides:
%   * the Madrid theme + book colour palette (matches the printed book)
%   * two book-style framing blocks:
%       \begin{bigpicture} ... \end{bigpicture}   ("the bigger picture")
%       \begin{underhood}   ... \end{underhood}    ("under the hood")
%     These mirror the book's `context` and `deepdive` tcolorboxes so the
%     same intuition-vs-mechanism scaffold the reader meets in the book
%     reappears on the slides.
%   * a Python listings style for code frames
%   * appendix conventions: \appendixstart drops a divider and switches
%     to non-numbered backup frames so "extra / less central" material
%     lives after the main talk without inflating the slide count.
%
% Usage:
%   \documentclass{beamer}
%   % =====================================================================
% finance-beamer.tex — shared Beamer style for the LLM-in-Finance course
% =====================================================================
% Reusable slide style. Each deck \input's this immediately after
% \documentclass{beamer}. It provides:
%   * the Madrid theme + book colour palette (matches the printed book)
%   * two book-style framing blocks:
%       \begin{bigpicture} ... \end{bigpicture}   ("the bigger picture")
%       \begin{underhood}   ... \end{underhood}    ("under the hood")
%     These mirror the book's `context` and `deepdive` tcolorboxes so the
%     same intuition-vs-mechanism scaffold the reader meets in the book
%     reappears on the slides.
%   * a Python listings style for code frames
%   * appendix conventions: \appendixstart drops a divider and switches
%     to non-numbered backup frames so "extra / less central" material
%     lives after the main talk without inflating the slide count.
%
% Usage:
%   \documentclass{beamer}
%   % =====================================================================
% finance-beamer.tex — shared Beamer style for the LLM-in-Finance course
% =====================================================================
% Reusable slide style. Each deck \input's this immediately after
% \documentclass{beamer}. It provides:
%   * the Madrid theme + book colour palette (matches the printed book)
%   * two book-style framing blocks:
%       \begin{bigpicture} ... \end{bigpicture}   ("the bigger picture")
%       \begin{underhood}   ... \end{underhood}    ("under the hood")
%     These mirror the book's `context` and `deepdive` tcolorboxes so the
%     same intuition-vs-mechanism scaffold the reader meets in the book
%     reappears on the slides.
%   * a Python listings style for code frames
%   * appendix conventions: \appendixstart drops a divider and switches
%     to non-numbered backup frames so "extra / less central" material
%     lives after the main talk without inflating the slide count.
%
% Usage:
%   \documentclass{beamer}
%   \input{../_style/finance-beamer.tex}
%   \title{...}\subtitle{...}\author{...}\institute{...}\date{\today}
%   \begin{document} ... \end{document}
% =====================================================================

\usetheme{Madrid}
\usecolortheme{whale}
\setbeamertemplate{navigation symbols}{}
\setbeamertemplate{caption}[numbered]
\setbeamercovered{transparent}

% --- Density controls: keep dense, equation-heavy frames on one slide ---
% Slightly smaller body text + tighter lists/blocks. Frame titles and the
% Madrid chrome keep their normal size, so the deck still reads as Madrid,
% just with more vertical room for content.
\setbeamerfont{normal text}{size=\small}
\setbeamertemplate{itemize/enumerate body begin}{\small}
\setbeamertemplate{itemize/enumerate subbody begin}{\footnotesize}
% Tighter block spacing.
\setbeamertemplate{block begin}{%
  \par\vskip2pt%
  \begin{beamercolorbox}[colsep*=.75ex,leftskip=1ex,rightskip=1ex]{block title}%
    \usebeamerfont*{block title}\insertblocktitle%
  \end{beamercolorbox}%
  {\parskip0pt\vskip-1pt}%
  \begin{beamercolorbox}[colsep*=.75ex,vmode,leftskip=1ex,rightskip=1ex]{block body}%
  \vskip1pt}
\setbeamertemplate{block end}{\end{beamercolorbox}\vskip2pt}

\usepackage{amsmath, amssymb}
\usepackage{booktabs}
\usepackage{xcolor}
\usepackage{listings}
\usepackage[skins, breakable]{tcolorbox}

% --- Book colour palette (kept in sync with book/preamble.tex) ---
\definecolor{linkblue}{RGB}{14, 75, 140}
\definecolor{deepdivebg}{RGB}{235, 244, 255}
\definecolor{narrativebg}{RGB}{255, 252, 240}
\definecolor{narrativerule}{RGB}{180, 130, 40}
\definecolor{steelblue}{RGB}{70, 130, 180}
\definecolor{forestgreen}{RGB}{34, 139, 34}
\definecolor{crimson}{RGB}{220, 20, 60}
\definecolor{darkorange}{RGB}{255, 140, 0}
\definecolor{codebg}{RGB}{248, 248, 248}
\definecolor{coderule}{RGB}{210, 210, 210}
\definecolor{keywordblue}{RGB}{0, 100, 180}
\definecolor{commentgray}{RGB}{120, 120, 120}
\definecolor{stringgreen}{RGB}{30, 130, 30}

% Tie the Madrid structural colour to the book's link blue.
\setbeamercolor{structure}{fg=linkblue}
\setbeamercolor{title}{fg=white}
\setbeamercolor{frametitle}{fg=white}

% --- "the bigger picture" : narrative / intuition framing -----------
\newtcolorbox{bigpicture}{%
  enhanced, breakable, boxsep=2pt,
  colback  = narrativebg,
  colframe = narrativerule,
  boxrule  = 0pt, toprule = 1.4pt, bottomrule = 0.4pt,
  leftrule = 0pt, rightrule = 0pt, arc = 0pt,
  left = 6pt, right = 6pt, top = 4pt, bottom = 5pt,
  fonttitle = \footnotesize\scshape\color{narrativerule},
  title = {the bigger picture}, attach title to upper = {\par\smallskip},
}

% --- "under the hood" : the mechanism / technical detail ------------
\newtcolorbox{underhood}{%
  enhanced, breakable, boxsep=2pt,
  colback  = deepdivebg,
  colframe = linkblue,
  boxrule  = 0pt, toprule = 1.4pt, bottomrule = 0.4pt,
  leftrule = 0pt, rightrule = 0pt, arc = 0pt,
  left = 6pt, right = 6pt, top = 4pt, bottom = 5pt,
  fonttitle = \footnotesize\scshape\color{linkblue},
  title = {under the hood}, attach title to upper = {\par\smallskip},
}

% --- Python code style ---------------------------------------------
\lstset{
  basicstyle    = \ttfamily\footnotesize,
  keywordstyle  = \color{keywordblue}\bfseries,
  commentstyle  = \color{commentgray}\itshape,
  stringstyle   = \color{stringgreen},
  showstringspaces = false,
  language      = Python,
  breaklines    = true,
  backgroundcolor = \color{codebg},
  frame         = single,
  rulecolor     = \color{coderule},
  numbers       = none,
  aboveskip     = 4pt, belowskip = 4pt,
}

% --- Appendix conventions ------------------------------------------
% \appendixstart{Title} : begins the "extra material" appendix. Frames
% after it are not counted toward the main deck's frame total, keeping
% the core talk lean while preserving the deeper / optional content.
\newcommand{\appendixstart}[1]{%
  \appendix
  \setbeamertemplate{frametitle continuation}{}%
  \begin{frame}[noframenumbering]
    \begin{center}
      {\Large\scshape\color{linkblue} Appendix}\\[4pt]
      {\large #1}\\[10pt]
      {\footnotesize Extra / optional material — beyond the core lecture.}
    \end{center}
  \end{frame}
}
% Use \begin{frame}[noframenumbering]{...} for each appendix frame.

%   \title{...}\subtitle{...}\author{...}\institute{...}\date{\today}
%   \begin{document} ... \end{document}
% =====================================================================

\usetheme{Madrid}
\usecolortheme{whale}
\setbeamertemplate{navigation symbols}{}
\setbeamertemplate{caption}[numbered]
\setbeamercovered{transparent}

% --- Density controls: keep dense, equation-heavy frames on one slide ---
% Slightly smaller body text + tighter lists/blocks. Frame titles and the
% Madrid chrome keep their normal size, so the deck still reads as Madrid,
% just with more vertical room for content.
\setbeamerfont{normal text}{size=\small}
\setbeamertemplate{itemize/enumerate body begin}{\small}
\setbeamertemplate{itemize/enumerate subbody begin}{\footnotesize}
% Tighter block spacing.
\setbeamertemplate{block begin}{%
  \par\vskip2pt%
  \begin{beamercolorbox}[colsep*=.75ex,leftskip=1ex,rightskip=1ex]{block title}%
    \usebeamerfont*{block title}\insertblocktitle%
  \end{beamercolorbox}%
  {\parskip0pt\vskip-1pt}%
  \begin{beamercolorbox}[colsep*=.75ex,vmode,leftskip=1ex,rightskip=1ex]{block body}%
  \vskip1pt}
\setbeamertemplate{block end}{\end{beamercolorbox}\vskip2pt}

\usepackage{amsmath, amssymb}
\usepackage{booktabs}
\usepackage{xcolor}
\usepackage{listings}
\usepackage[skins, breakable]{tcolorbox}

% --- Book colour palette (kept in sync with book/preamble.tex) ---
\definecolor{linkblue}{RGB}{14, 75, 140}
\definecolor{deepdivebg}{RGB}{235, 244, 255}
\definecolor{narrativebg}{RGB}{255, 252, 240}
\definecolor{narrativerule}{RGB}{180, 130, 40}
\definecolor{steelblue}{RGB}{70, 130, 180}
\definecolor{forestgreen}{RGB}{34, 139, 34}
\definecolor{crimson}{RGB}{220, 20, 60}
\definecolor{darkorange}{RGB}{255, 140, 0}
\definecolor{codebg}{RGB}{248, 248, 248}
\definecolor{coderule}{RGB}{210, 210, 210}
\definecolor{keywordblue}{RGB}{0, 100, 180}
\definecolor{commentgray}{RGB}{120, 120, 120}
\definecolor{stringgreen}{RGB}{30, 130, 30}

% Tie the Madrid structural colour to the book's link blue.
\setbeamercolor{structure}{fg=linkblue}
\setbeamercolor{title}{fg=white}
\setbeamercolor{frametitle}{fg=white}

% --- "the bigger picture" : narrative / intuition framing -----------
\newtcolorbox{bigpicture}{%
  enhanced, breakable, boxsep=2pt,
  colback  = narrativebg,
  colframe = narrativerule,
  boxrule  = 0pt, toprule = 1.4pt, bottomrule = 0.4pt,
  leftrule = 0pt, rightrule = 0pt, arc = 0pt,
  left = 6pt, right = 6pt, top = 4pt, bottom = 5pt,
  fonttitle = \footnotesize\scshape\color{narrativerule},
  title = {the bigger picture}, attach title to upper = {\par\smallskip},
}

% --- "under the hood" : the mechanism / technical detail ------------
\newtcolorbox{underhood}{%
  enhanced, breakable, boxsep=2pt,
  colback  = deepdivebg,
  colframe = linkblue,
  boxrule  = 0pt, toprule = 1.4pt, bottomrule = 0.4pt,
  leftrule = 0pt, rightrule = 0pt, arc = 0pt,
  left = 6pt, right = 6pt, top = 4pt, bottom = 5pt,
  fonttitle = \footnotesize\scshape\color{linkblue},
  title = {under the hood}, attach title to upper = {\par\smallskip},
}

% --- Python code style ---------------------------------------------
\lstset{
  basicstyle    = \ttfamily\footnotesize,
  keywordstyle  = \color{keywordblue}\bfseries,
  commentstyle  = \color{commentgray}\itshape,
  stringstyle   = \color{stringgreen},
  showstringspaces = false,
  language      = Python,
  breaklines    = true,
  backgroundcolor = \color{codebg},
  frame         = single,
  rulecolor     = \color{coderule},
  numbers       = none,
  aboveskip     = 4pt, belowskip = 4pt,
}

% --- Appendix conventions ------------------------------------------
% \appendixstart{Title} : begins the "extra material" appendix. Frames
% after it are not counted toward the main deck's frame total, keeping
% the core talk lean while preserving the deeper / optional content.
\newcommand{\appendixstart}[1]{%
  \appendix
  \setbeamertemplate{frametitle continuation}{}%
  \begin{frame}[noframenumbering]
    \begin{center}
      {\Large\scshape\color{linkblue} Appendix}\\[4pt]
      {\large #1}\\[10pt]
      {\footnotesize Extra / optional material — beyond the core lecture.}
    \end{center}
  \end{frame}
}
% Use \begin{frame}[noframenumbering]{...} for each appendix frame.

%   \title{...}\subtitle{...}\author{...}\institute{...}\date{\today}
%   \begin{document} ... \end{document}
% =====================================================================

\usetheme{Madrid}
\usecolortheme{whale}
\setbeamertemplate{navigation symbols}{}
\setbeamertemplate{caption}[numbered]
\setbeamercovered{transparent}

% --- Density controls: keep dense, equation-heavy frames on one slide ---
% Slightly smaller body text + tighter lists/blocks. Frame titles and the
% Madrid chrome keep their normal size, so the deck still reads as Madrid,
% just with more vertical room for content.
\setbeamerfont{normal text}{size=\small}
\setbeamertemplate{itemize/enumerate body begin}{\small}
\setbeamertemplate{itemize/enumerate subbody begin}{\footnotesize}
% Tighter block spacing.
\setbeamertemplate{block begin}{%
  \par\vskip2pt%
  \begin{beamercolorbox}[colsep*=.75ex,leftskip=1ex,rightskip=1ex]{block title}%
    \usebeamerfont*{block title}\insertblocktitle%
  \end{beamercolorbox}%
  {\parskip0pt\vskip-1pt}%
  \begin{beamercolorbox}[colsep*=.75ex,vmode,leftskip=1ex,rightskip=1ex]{block body}%
  \vskip1pt}
\setbeamertemplate{block end}{\end{beamercolorbox}\vskip2pt}

\usepackage{amsmath, amssymb}
\usepackage{booktabs}
\usepackage{xcolor}
\usepackage{listings}
\usepackage[skins, breakable]{tcolorbox}

% --- Book colour palette (kept in sync with book/preamble.tex) ---
\definecolor{linkblue}{RGB}{14, 75, 140}
\definecolor{deepdivebg}{RGB}{235, 244, 255}
\definecolor{narrativebg}{RGB}{255, 252, 240}
\definecolor{narrativerule}{RGB}{180, 130, 40}
\definecolor{steelblue}{RGB}{70, 130, 180}
\definecolor{forestgreen}{RGB}{34, 139, 34}
\definecolor{crimson}{RGB}{220, 20, 60}
\definecolor{darkorange}{RGB}{255, 140, 0}
\definecolor{codebg}{RGB}{248, 248, 248}
\definecolor{coderule}{RGB}{210, 210, 210}
\definecolor{keywordblue}{RGB}{0, 100, 180}
\definecolor{commentgray}{RGB}{120, 120, 120}
\definecolor{stringgreen}{RGB}{30, 130, 30}

% Tie the Madrid structural colour to the book's link blue.
\setbeamercolor{structure}{fg=linkblue}
\setbeamercolor{title}{fg=white}
\setbeamercolor{frametitle}{fg=white}

% --- "the bigger picture" : narrative / intuition framing -----------
\newtcolorbox{bigpicture}{%
  enhanced, breakable, boxsep=2pt,
  colback  = narrativebg,
  colframe = narrativerule,
  boxrule  = 0pt, toprule = 1.4pt, bottomrule = 0.4pt,
  leftrule = 0pt, rightrule = 0pt, arc = 0pt,
  left = 6pt, right = 6pt, top = 4pt, bottom = 5pt,
  fonttitle = \footnotesize\scshape\color{narrativerule},
  title = {the bigger picture}, attach title to upper = {\par\smallskip},
}

% --- "under the hood" : the mechanism / technical detail ------------
\newtcolorbox{underhood}{%
  enhanced, breakable, boxsep=2pt,
  colback  = deepdivebg,
  colframe = linkblue,
  boxrule  = 0pt, toprule = 1.4pt, bottomrule = 0.4pt,
  leftrule = 0pt, rightrule = 0pt, arc = 0pt,
  left = 6pt, right = 6pt, top = 4pt, bottom = 5pt,
  fonttitle = \footnotesize\scshape\color{linkblue},
  title = {under the hood}, attach title to upper = {\par\smallskip},
}

% --- Python code style ---------------------------------------------
\lstset{
  basicstyle    = \ttfamily\footnotesize,
  keywordstyle  = \color{keywordblue}\bfseries,
  commentstyle  = \color{commentgray}\itshape,
  stringstyle   = \color{stringgreen},
  showstringspaces = false,
  language      = Python,
  breaklines    = true,
  backgroundcolor = \color{codebg},
  frame         = single,
  rulecolor     = \color{coderule},
  numbers       = none,
  aboveskip     = 4pt, belowskip = 4pt,
}

% --- Appendix conventions ------------------------------------------
% \appendixstart{Title} : begins the "extra material" appendix. Frames
% after it are not counted toward the main deck's frame total, keeping
% the core talk lean while preserving the deeper / optional content.
\newcommand{\appendixstart}[1]{%
  \appendix
  \setbeamertemplate{frametitle continuation}{}%
  \begin{frame}[noframenumbering]
    \begin{center}
      {\Large\scshape\color{linkblue} Appendix}\\[4pt]
      {\large #1}\\[10pt]
      {\footnotesize Extra / optional material — beyond the core lecture.}
    \end{center}
  \end{frame}
}
% Use \begin{frame}[noframenumbering]{...} for each appendix frame.


\title{Practical Session 16: AI/ML Foundations for Financial Text}
\subtitle{Hands-On: Lexicon Sentiment, BoW vs.\ Embeddings, and a Mini Attention Trace}
\author{Juan F.\ Imbet}
\institute{Paris Dauphine -- PSL University}
\date{\today}

\begin{document}

\begin{frame}\titlepage\end{frame}

% =====================================================================
% FRAME: Session Overview
% =====================================================================
\begin{frame}{Session Overview}

\textbf{Duration:} 1 hour (50 min activities + 10 min debrief)

\medskip
\begin{enumerate}
  \item \textbf{Recap} (5 min): the AI/ML/DL/LLM hierarchy; the attention formula.
  \item \textbf{Problem 1} (12 min): Loughran--McDonald dictionary sentiment on a
        filing snippet --- by hand --- and why a \emph{general} lexicon misleads.
  \item \textbf{Problem 2} (10 min): Show that mean BoW pooling is order-blind on
        a financial sentence; quantify the failure.
  \item \textbf{Problem 3} (8 min): Trace a $3\times3$ self-attention step on a
        financial sequence (uses the recap formula).
  \item \textbf{Case Study} (15 min): Python --- dictionary vs.\ embedding sentiment on
        SEC snippets (\texttt{code/notebooks/16-ai-ml-finance-text/}).
\end{enumerate}

\medskip
\begin{block}{Discussion (10 min)}
When is a 1990s-vintage dictionary still the \emph{right} tool, and when must you
reach for an embedding or an LLM?
\end{block}
\end{frame}

% =====================================================================
% FRAME: Recap 1
% =====================================================================
\begin{frame}{Recap: The Hierarchy and Where Text Fits}

\begin{bigpicture}
AI $\supset$ ML $\supset$ DL $\supset$ LLM. \textbf{Statistical} AI (learn
parameters from data) dominates finance over \textbf{symbolic} AI (hand-coded
rules) because markets are data-rich --- but symbolic rules return for compliance
and interpretability.
\end{bigpicture}

\medskip
\textbf{Why text is hard (carry into the problems):}
\begin{itemize}
  \item \textbf{Ambiguous:} $\sim$3/4 of general ``negative'' words are
        neutral/positive in filings $\Rightarrow$ use a \emph{finance} lexicon.
  \item \textbf{Sparse:} 50k-word vocabulary $\gg$ observations.
  \item \textbf{Strategic:} disclosures written by lawyers --- not unbiased signals.
\end{itemize}
\end{frame}

% =====================================================================
% FRAME: Recap 2
% =====================================================================
\begin{frame}{Recap: Two Formulas You Will Use}

\textbf{Dictionary sentiment} of a document:
\[
  \mathrm{Sentiment} = \frac{\#\text{positive words} - \#\text{negative words}}
                            {\#\text{total words}}.
\]

\medskip
\textbf{Mean bag-of-words embedding} (order-blind by construction):
\[
  \bar{\varphi}(d) = \frac{1}{n}\sum_{i=1}^{n}\varphi(w_i).
\]

\medskip
\textbf{Scaled dot-product attention} (Problem 3):
\[
  \mathrm{Attention}(Q,K,V) = \mathrm{softmax}\!\left(\frac{QK^\top}{\sqrt{d_k}}\right)V,
  \qquad
  \mathrm{softmax}(\mathbf{x})_i = \frac{e^{x_i}}{\sum_j e^{x_j}}.
\]

\medskip
\emph{Tip:} subtract $\max(\mathbf{x})$ before exponentiating for numerical stability.
\end{frame}

% =====================================================================
% FRAME: Problem 1
% =====================================================================
\begin{frame}{Problem 1: Dictionary Sentiment, and Why Domain Matters}

\textbf{Snippet} (10-K Risk Factors, stylised):
\begin{quote}\small
``Demand remained \textbf{positive}, but \textbf{litigation} exposure and
macroeconomic \textbf{uncertainty} could adversely affect results; management
believes the outlook is \textbf{strong}.''
\end{quote}

\medskip
\textbf{Task A.} Using Loughran--McDonald's 6 categories
(negative, positive, uncertainty, litigious, strong/weak modal), classify the
four bold words. Which category does ``litigation'' / ``uncertainty'' fall into?

\medskip
\textbf{Task B.} A \emph{general} (Harvard Inquirer) dictionary tags
``litigation,'' ``liability,'' and ``tax'' as negative. The chapter notes
$\sim$3/4 of Harvard ``negative'' words are neutral/positive in filings.
What does that imply for a naive negative-word count on legal-heavy filings?

\medskip
\textbf{Task C.} Compute the net sentiment $(\#\text{pos}-\#\text{neg})/\#\text{total}$
for the snippet (10 words after dropping stop-words), using a finance lexicon
where only \emph{negative} \& \emph{positive} categories count toward polarity.
\end{frame}

% =====================================================================
% FRAME: Problem 2
% =====================================================================
\begin{frame}{Problem 2: Mean Pooling is Order-Blind}

\textbf{Two sentences with opposite meaning:}
\begin{itemize}
  \item $A$: ``Revenue rose, costs fell.''
  \item $B$: ``Costs rose, revenue fell.''
\end{itemize}

\medskip
\textbf{Task A.} The two sentences contain the \emph{same multiset of words}.
Show that the mean BoW embedding $\bar\varphi(A)=\bar\varphi(B)$ exactly. Why?

\medskip
\textbf{Task B.} A long-only equity signal goes \emph{long} on positive-sentiment
filings. If your representation cannot distinguish $A$ from $B$, what failure mode
appears on filings where the \emph{subject} of ``rose'' vs.\ ``fell'' matters?

\medskip
\textbf{Task C.} Name the architectural ingredient (from the lecture) that
restores order sensitivity, and state in one line \emph{why} it can connect
``revenue'' to ``rose'' across a long sentence where an RNN struggles.
\end{frame}

% =====================================================================
% FRAME: Problem 3
% =====================================================================
\begin{frame}{Problem 3: A $3\times3$ Self-Attention Trace}

\textbf{Sequence:} \texttt{"risk"}, \texttt{"increased"}, \texttt{"materially"}.
\quad Head dimension $d_k = 2$.

\medskip
\[
Q = \begin{pmatrix} 1 & 0 \\ 1 & 1 \\ 0 & 1 \end{pmatrix},\quad
K = \begin{pmatrix} 1 & 0 \\ 1 & 1 \\ 0 & 1 \end{pmatrix},\quad
V = \begin{pmatrix} 2 & 0 \\ 0 & 2 \\ 1 & 1 \end{pmatrix}
\]

\medskip
\textbf{Task A.} Compute $S = QK^\top$ (entry $(i,j)=\mathbf{q}_i^\top\mathbf{k}_j$).

\textbf{Task B.} Scale: $\tilde S = S/\sqrt{2}$ (two decimals).

\textbf{Task C.} Softmax row-wise $\to A$; \quad
\textbf{Task D.} output $O = AV$.

\medskip
\textbf{Discussion:} Which token does \texttt{"risk"} attend to most? If we set
$d_k=64$ with the same raw dot products, are the weights \emph{sharper} or flatter?
\end{frame}

% =====================================================================
% FRAME: Case Study Intro
% =====================================================================
\begin{frame}{Case Study: Dictionary vs.\ Embedding Sentiment}

\textbf{Goal:} on a handful of SEC-style snippets, compare a
\textbf{Loughran--McDonald-style dictionary} score against a \textbf{transformer
embedding} similarity score, and see where they disagree.

\medskip
\textbf{Stack:}
\begin{itemize}
  \item \texttt{sentence\_transformers} --- contextual embeddings (SBERT).
  \item A small finance positive/negative word list --- the dictionary baseline.
\end{itemize}

\medskip
\textbf{Notebook:} \texttt{code/notebooks/16-ai-ml-finance-text/demo.ipynb}

\medskip
\begin{alertblock}{Setup}
\texttt{pip install sentence-transformers}\\
No API key needed --- embeddings run locally on CPU.
\end{alertblock}
\end{frame}

% =====================================================================
% FRAME: Case Study Code
% =====================================================================
\begin{frame}[fragile]{Case Study: The Code (1/2) --- Data \& Dictionary}

\begin{lstlisting}
from sentence_transformers import SentenceTransformer, util

SNIPPETS = [
    "Revenue rose 12% and margins expanded.",          # clearly positive
    "Litigation exposure and uncertainty increased.",  # negative, legal-heavy
    "Costs fell, revenue fell, outlook unchanged.",     # order-sensitive
]
POS = {"rose", "expanded", "growth", "strong", "gain"}
NEG = {"fell", "litigation", "uncertainty", "loss", "decline"}

def dict_sentiment(text):
    w = text.lower().replace(",", " ").replace(".", " ").split()
    return (sum(t in POS for t in w) - sum(t in NEG for t in w)) / len(w)
\end{lstlisting}
\end{frame}

% =====================================================================
% FRAME: Case Study Code (2/2)
% =====================================================================
\begin{frame}[fragile]{Case Study: The Code (2/2) --- Embedding Similarity}

\begin{lstlisting}
model = SentenceTransformer("all-MiniLM-L6-v2")
anchor = model.encode("This is good news for shareholders.",
                      normalize_embeddings=True)
for s in SNIPPETS:
    d = dict_sentiment(s)
    emb = model.encode(s, normalize_embeddings=True)
    sim = float(util.cos_sim(emb, anchor))
    print(f"dict={d:+.3f}  embed_sim={sim:+.3f}  | {s}")
\end{lstlisting}
\end{frame}

% =====================================================================
% FRAME: Case Study Questions
% =====================================================================
\begin{frame}{Case Study: Diagnostic Questions}

\textbf{Q1.} For each snippet, report the dictionary score and the embedding
similarity. Where do the two \emph{disagree} in sign or ranking?

\medskip
\textbf{Q2.} Snippet 3 (``Costs fell, revenue fell\dots'') contains the
\emph{negative} word ``fell'' twice. Does the dictionary call it negative? Is that
correct, given costs falling is \emph{good}? What does this reveal about
\textbf{context}-free scoring?

\medskip
\textbf{Q3.} Replace ``uncertainty'' with ``opportunity'' in snippet 2. Does the
dictionary score flip? Does the embedding similarity move \emph{more} or
\emph{less} than the dictionary? Why?

\medskip
\begin{block}{Optional extension}
Add ``positive'' once to snippet 2 (recall ``positive'' is \emph{cautious} in
credit-rating language). Does the dictionary over-credit it? Sketch how you would
disambiguate using the document's genre.
\end{block}
\end{frame}

% =====================================================================
% FRAME: Discussion
% =====================================================================
\begin{frame}{Discussion: Dictionary vs.\ Embedding vs.\ LLM}

For each scenario, pick \textbf{dictionary}, \textbf{embedding/encoder}, or
\textbf{LLM}, and justify on cost, auditability, label availability, and context.

\medskip
\begin{enumerate}
  \item A regulator must \emph{reproduce} a sentiment score from a published,
        auditable rule on 200k filings --- transparency is mandatory.
  \item A desk wants to rank 50k earnings-call paragraphs by similarity to a
        ``supply-chain risk'' theme, with no labeled training data.
  \item An analyst asks: ``Summarize how this 10-K's risk language changed
        versus last year, and flag any new litigation exposure.''
\end{enumerate}

\medskip
\begin{block}{Frame your answer around}
Transparency / reproducibility, context sensitivity, labeled-data needs,
generation vs.\ scoring, cost, and the strategic production of the text.
\end{block}
\end{frame}

% =====================================================================
% FRAME: Solution 1
% =====================================================================
\begin{frame}{Solution: Problem 1 --- Dictionary Sentiment}

\textbf{Task A --- categories:}
``positive'' $\to$ \emph{positive}; ``strong'' $\to$ \emph{strong modal} (and
positive-leaning); ``litigation'' $\to$ \emph{litigious}; ``uncertainty'' $\to$
\emph{uncertainty}. Note: litigious and uncertainty are \emph{not} polarity-negative
per se --- they are their own categories.

\medskip
\textbf{Task B:} A general dictionary tags ``litigation'' as negative, so a naive
negative count \alert{systematically over-penalises legal-heavy filings} ---
precisely the Loughran--McDonald finding ($\sim$3/4 of Harvard negatives are
neutral/positive in filings).

\medskip
\textbf{Task C:} With only \emph{positive}/\emph{negative} counting toward polarity:
positives $=\{$positive, strong$\}=2$, negatives $=0$ (litigious/uncertainty
excluded). Net $=(2-0)/10 = \mathbf{+0.20}$. A general dictionary would instead
score this \emph{negative} --- the wrong sign.
\end{frame}

% =====================================================================
% FRAME: Solution 2
% =====================================================================
\begin{frame}{Solution: Problem 2 --- Order-Blindness}

\textbf{Task A:} Mean pooling sums per-word embeddings then divides by $n$.
Addition is commutative, so any permutation of the \emph{same} words yields the
\emph{identical} mean. $A$ and $B$ share the multiset
$\{$revenue, rose, costs, fell$\}$ $\Rightarrow \bar\varphi(A)=\bar\varphi(B)$ exactly.

\medskip
\textbf{Task B:} A sentiment-driven long-only signal cannot tell ``revenue rose''
(bullish) from ``revenue fell'' (bearish): it would take the \emph{same} position
on opposite fundamentals --- a sign error that no amount of training fixes if the
representation discards order.

\medskip
\textbf{Task C:} \textbf{Self-attention} restores order sensitivity (with
positional information). It can directly connect ``revenue'' to ``rose'' because
every position attends to every other in one step --- no sequential path to
decay, unlike an RNN whose gradient vanishes over long spans.
\end{frame}

% =====================================================================
% FRAME: Solution 3
% =====================================================================
\begin{frame}{Solution: Problem 3 --- Attention Trace}

Here $K=Q$, so $S=QQ^\top$:
\[
S = \begin{pmatrix} 1 & 1 & 0 \\ 1 & 2 & 1 \\ 0 & 1 & 1 \end{pmatrix},
\qquad
\tilde S = S/\sqrt{2} \approx
\begin{pmatrix} 0.71 & 0.71 & 0.00 \\ 0.71 & 1.41 & 0.71 \\ 0.00 & 0.71 & 0.71 \end{pmatrix}.
\]

\textbf{Row-wise softmax} $\to A$; \quad \textbf{output} $O = AV$ with
$V=\big[(2,0),(0,2),(1,1)\big]^\top$:
\[
A \approx
\begin{pmatrix} 0.422 & 0.422 & 0.156 \\ 0.270 & 0.461 & 0.270 \\ 0.156 & 0.422 & 0.422 \end{pmatrix},
\qquad
O \approx
\begin{pmatrix} 1.00 & 1.00 \\ 0.81 & 1.19 \\ 0.73 & 1.27 \end{pmatrix}.
\]

\textbf{Discussion:} \texttt{"risk"} attends most to itself and \texttt{"increased"}
(tied at $0.422$). At $d_k=64$ the same raw dot products are divided by a larger
$\sqrt{d_k}$ \emph{only if unscaled} --- but with the same \emph{pre-scaling}
magnitudes the larger logits make softmax \alert{sharper}; the $\sqrt{d_k}$ scaling
exists precisely to counter that saturation.
\end{frame}

% =====================================================================
% FRAME: Wrap-up
% =====================================================================
\begin{frame}{Wrap-Up: Key Takeaways}

\textbf{From Problem 1:} a \emph{domain} lexicon (Loughran--McDonald) avoids the
systematic negative bias a general dictionary inflicts on legal-heavy filings ---
the canonical motivation for finance-specific NLP.

\medskip
\textbf{From Problem 2:} bag-of-words mean pooling is provably order-blind;
opposite-meaning sentences collapse to one vector. Order matters in finance, and
self-attention is the fix.

\medskip
\textbf{From Problem 3:} attention is a weighted average of value vectors; weights
are softmaxed scaled dot products; $\sqrt{d_k}$ keeps the softmax out of its
saturated, vanishing-gradient regime.

\medskip
\textbf{From the Case Study:} dictionary and embedding scores disagree exactly
where \emph{context} matters (``costs fell'' is good; ``positive'' is cautious in
ratings) --- the case for moving up the ladder to encoders and LLMs.

\medskip
\begin{block}{Next practical}
Building the transformer block end-to-end (LLM-foundations): multi-head
attention, positional encoding, and a first fine-tuning pass.
\end{block}
\end{frame}

% =====================================================================
\appendixstart{Stretch Problems}
% =====================================================================

\begin{frame}[noframenumbering]{Stretch Problem A: tf-idf by Hand}

Three one-sentence ``documents'':
\begin{itemize}
  \item $d_1$: ``revenue revenue growth''
  \item $d_2$: ``revenue decline''
  \item $d_3$: ``growth growth growth''
\end{itemize}

\medskip
\textbf{Task.} With $\mathrm{tf}$ = raw count and
$\mathrm{idf}(w)=\log\frac{N}{df(w)}$ ($N=3$), compute the tf-idf weight of
``revenue'' in $d_1$ and of ``growth'' in $d_3$. Which term is up-weighted, and why
does tf-idf \emph{down-weight} a word appearing in every document?

\medskip
\begin{block}{Point}
tf-idf rewards rare, document-specific words --- the BoW refinement that made
1990s filing analysis workable before embeddings existed.
\end{block}
\end{frame}

\begin{frame}[noframenumbering]{Stretch Problem B: Reflexivity in a Backtest}

A text-based ``negative-sentiment $\Rightarrow$ short'' signal backtests with a
Sharpe of 1.4 over 2015--2019.

\medskip
\textbf{Task A.} Name two reasons the live Sharpe may be far lower: (i) the signal
becomes \emph{public}; (ii) the strategy's own trades move prices (reflexivity).

\medskip
\textbf{Task B.} The chapter warns LLMs can \emph{reconstruct historical series
from parametric memory}, producing spurious predictability. If your sentiment
model was pre-trained on data \emph{overlapping} the backtest window, what bias
contaminates the result, and how would you guard against it?

\medskip
\begin{block}{Point}
AI systems in finance are \emph{actors within} the market, not observers ---
look-ahead and reflexivity are first-order design concerns, not footnotes.
\end{block}
\end{frame}

\end{document}
